% =========================================================
% SECCIÓN 3.1 - INTRODUCCIÓN AL CAPÍTULO 3
% =========================================================

El Capítulo \ref{chap:instrumentacion} estableció la base instrumental y metrológica que garantiza la trazabilidad del dato crudo proveniente de los reactores batch de bioconversión. Mediante sensores NDIR de bajo costo, calibración frente a patrones certificados y compensación cruzada por temperatura y humedad relativa, se obtuvo una serie temporal de flujos másicos de \ch{CO2} —expresada como $\dot{m}_{\mathrm{CO_2}}(t)$ en g/min— con una incertidumbre expandida del $\pm 12\%$ y un MAPE del $9.8\%$ respecto al instrumento de referencia. Estos datos validados satisfacen la hipótesis H1 ($R > 0.85$, MAPE $< 15\%$) y constituyen la entrada empírica indispensable para cerrar el ciclo del sistema híbrido: sensores $\rightarrow$ modelo $\rightarrow$ estimador $\rightarrow$ métricas ambientales.

No obstante, la señal $\dot{m}_{\mathrm{CO_2}}(t)$, por sí sola, permanece como un \emph{proxy} físico cuya interpretación bioenergética no es inmediata. Un flujo de \ch{CO2} de, por ejemplo, \num{0.45} g/min en el día 7 del ciclo larval puede reflejar una tasa de asimilación elevada con metabolismo aeróbico eficiente, o bien una fracción creciente de respiración microbiana anaerobia compensada parcialmente, condición que el sensor de \ch{CH4} (Capítulo \ref{chap:metricas}) detectaría como evento de anaerobiosis. La necesidad de \emph{desagregar} el flujo gaseoso en sus componentes fisiológicos —mantenimiento somático, síntesis de biomasa estructural, acumulación de lípidos de reserva y respiración microbiana asociada— exige un modelo mecanicista que articule la cinética de asimilación del sustrato con la termodinámica del crecimiento larval.

En este contexto, el presente capítulo desarrolla la generalización del modelo Dynamic Energy Budget (DEB) propuesto por \parencite{eriksenDynamicModellingFeed2022} desde la escala individual larval hasta la escala de reactor batch. El objetivo central es formular un sistema de ecuaciones diferenciales ordinarias (EDO) que prediga la tasa de producción de \ch{CO2} como resultado emergente de los flujos metabólicos de carbono: ingesta, asimilación, distribución entre compartimentos de reserva y estructura, y costos energéticos de mantenimiento y síntesis. El modelo DEB se parametriza con semillas tomadas de la literatura (Tabla S1 de Eriksen) y deja variables libres —$\tau_h$, $\beta_0$, $\rho_{\mathrm{conv}}$, $T_A$, $T_{\mathrm{opt}}$, $T_H$— para su calibración local en el Capítulo \ref{chap:hibrido} mediante confrontación con las series de sensores NDIR.

La estructura del capítulo se organiza en ocho secciones. La Sección \ref{sec:cap3_fundamentos_deb} presenta los fundamentos del modelo DEB para \emph{H. illucens}, incluyendo la arquitectura de compartimentos de carbono y los flujos metabólicos que generan \ch{CO2}. La Sección \ref{sec:cap3_variables} define las variables de estado, sus dominios físicos y las condiciones iniciales derivadas del protocolo experimental del Capítulo \ref{chap:instrumentacion}. La Sección \ref{sec:cap3_formulacion_ode} formula el sistema de EDO que gobierna la dinámica del sustrato, los asimilados, la biomasa estructural, los lípidos de reserva y la acumulación de \ch{CO2}. La Sección \ref{sec:cap3_cinetica} introduce la cinética de asimilación tipo Holling II y la modulación térmica mediante el modelo de Sharpe--Schoolfield. La Sección \ref{sec:cap3_parametros} detalla los parámetros del modelo, sus unidades, rangos semilla y las restricciones físicas que garantizan la disipación positiva y la conservación de masa. La Sección \ref{sec:cap3_analisis} analiza la consistencia dimensional, el balance de carbono y el comportamiento cualitativo del sistema; las demostraciones formales de existencia, unicidad y estabilidad de las soluciones se remiten al Anexo \ref{anexo:formulacion_matematica} para no interrumpir el hilo argumental del capítulo. Finalmente, la Sección \ref{sec:cap3_conclusiones} sintetiza los hallazgos y establece el puente hacia la integración híbrida con el estimador de biomasa del Capítulo \ref{chap:hibrido}.

Es importante precisar el alcance del capítulo. El metano (\ch{CH4}) no se incorpora como variable de estado en las EDE; dado que el proceso opera predominantemente en régimen aeróbico, el \ch{CH4} se trata como residual diagnóstico de eventos de anaerobiosis transitoria, cuya detección y cuantificación se abordan en el Capítulo \ref{chap:metricas}. Asimismo, el modelo asume homogeneidad espacial en el reactor batch, de modo que los gradientes de oxigenación y temperatura dentro del lecho de sustrato se promedian en una única variable ambiental efectiva. Esta simplificación, aunque conservadora, es consistente con la escala de los reactores de \num{20} L utilizados en la caracterización instrumental y permite mantener el sistema en un número manejable de ecuaciones sin pérdida sustancial de capacidad predictiva a escala piloto \parencite{padmanabhaComprehensiveDynamicGrowth2020,rossiEstimatingDynamicsGreenhouse2024}.

En síntesis, este capítulo transforma los datos crudos validados del Capítulo \ref{chap:instrumentacion} en una narrativa bioenergética cuantitativa: interpreta \emph{qué} procesos biológicos producen el \ch{CO2} medido, \emph{cuánto} carbono se destina a cada compartimento y \emph{bajo qué} condiciones ambientales la predicción mecanicista es fiable. El resultado es un modelo DEB escalado a reactor batch, listo para ser calibrado y validado contra sensores en el Capítulo \ref{chap:hibrido}.
